% 原书第 101—102 页；第 7.3 节全文。
\section{讨论}
\label{sec:ch07-discussion}

本章证明了：如果使用算法库中包含 HAL 估计量的超级学习器作为初始估计量，那么，特定处理下均值的一步 TMLE 就具有渐近有效性，因而平均处理效应（average treatment effect）的估计也具有渐近有效性。

\begin{sourcebox}
关键在于，这些干扰参数的 HAL 估计量，能够以快于临界速率 $n^{-1/4}$ 的速率收敛到各自的真值。我们在证明这一点时，只假设了强正值性、损失函数一致有界，以及干扰参数的变差范数以某个通用常数为界。
\end{sourcebox}

还可以仅假设真实干扰参数具有有限的变差范数，从而证明 TMLE 的渐近有效性。这就允许模型关于变差范数无界，但仍然假设模型关于上确界范数（supremum norm）具有一致的界，使损失函数和有效影响曲线一致有界。在这种情况下，使用一个由有界模型 $\mathcal M_n$ 组成的筛（sieve），允许 $\mathcal M_n$ 的通用界随 $n$ 增大或减小，从而当 $n$ 足够大时，真实干扰参数便落在相应于 $n$ 的模型 $\mathcal M_n$ 中；即当 $n>N_0=N(P_0)$ 时，有 $P_0\in\mathcal M_n$。使用交叉验证 TMLE（cross-validated TMLE，CV-TMLE）时，甚至不需要通过筛来控制变差范数界。对于 TMLE 和 CV-TMLE，我们甚至可以允许筛 $\mathcal M_n$ 的上确界范数界趋于无穷，只要其增长足够缓慢，使超级学习器的预言不等式（oracle inequality）中的二阶项仍然以快于 $n^{-1/2}$ 的速率收敛到零。我们还展示了，这些结果可以直接推广到一般的模型和目标参数。关于一般模型以及特定处理下均值这一例子的精确定理，请参阅 van der Laan（2017）。

由于我们对干扰参数 $Q$ 和 $G$ 使用了 HAL 估计量，就不需要依赖有效影响曲线的双重稳健性（double robustness），即当 $Q=Q_0$ 或 $G=G_0$ 中任一个成立时，均有
$R_{20}((Q,G),(Q_0,G_0))=0$。
在 van der Laan（2014b）中，我们展示了：对于这样的双重稳健估计问题，可以构造一个 TMLE，即使两个干扰参数估计量中有一个不一致，它仍保持渐近线性。例如，如果 $Q_n$ 收敛到一个设定错误的 $Q$，那么，余项
$R_{20}((Q_n,G_n),(Q_0,G_0))$
就不再是二阶项，而是表现为
$R_{20}((Q,G_n),(Q_0,G_0))$。
后者可以写为 $G_n$ 的某个函数减去同一函数在 $G_0$ 处的值。因此，为了使
$R_{20}((Q,G_n),(Q_0,G_0))$
具有渐近线性，$G_n$ 必须是一个 TMLE，靶向于
$G\mapsto R_{20}((Q,G),(Q_0,G_0))$
中所需的泛函。事实上，van der Laan（2014b）提出的 TMLE 同时包含 $Q$ 和 $G$ 的扰动模型，使得所得 $(Q_0,G_0)$ 的 TMLE $(Q_n^*,G_n^*)$，不仅靶向于目标参数 $\psi_0$，也靶向于这些额外的干扰参数。尽管这些更复杂的做法不会改善 HAL-TMLE 的渐近表现，它们仍可能改善 HAL-TMLE 的有限样本表现。

